\section{Influence of ionic strength on particle dielectrophoretic trapping}
\subsubsection{Abstract}
Showcasing that our setup can detect and actuate proteins, giving information about their physical properties. Using Encapsulin we test ideal concentrations for trapping of proteins in water, then we change the ionic strength and see if the trapping capabilities are still the same. We compare this with particles of the same site.
Protein embedded dipole moment could make proteins in high ionic strength trappable, but particles not. 
	%	Situation: \\
	%	Complication: \\
	%	Question: \\
	%	Answer: \\
	%	Visualisation: \\
    Situation: \\
	Complication: \\
	Question: \\
	Answer: \\
	Visualisation: \\
    \\\\
    Figure list:
    \begin{description}
        \item [Supplementary figure] encasulin production and purity
        \item [Figure 1] ionic strength/conductivity vs trapping prop particles and DNA
        \item [Figure 3] ionic strength/conductivity vs trapping probability Encapsulin
        \item [Figure 2] ionic strength vs duration and intensity distribution for particles DNA DNA
        \item [Figure 4] ionic strength vs duration and intensity distribution for proteins
    \end{description}
\subsubsection{Introduction}
    Situation: We want to investigate protein dipole-moments\\
	Complication: no technology exists that is able to do that yet\\
	Question: how can we measure protein dipole moment? \\
	Answer: using dep\\
	Visualisation: refer to introduction with proteinDep figure\\
    \\\\How
    Situation: protein dep requires an electric field\\
	Complication: proteins require biological "high" salt conditions\\
	Question: Can we still trap in high salt conditions?\\
	Answer: We tested it here, investigating trapping probability at various ionic strengths. Proteins might be trappable, DNA definitely is, but particles might not. \\
	Visualisation: not needed \\
    \\\\
    Situation: We want to test ionic strength on proteins\\
	Complication: We need a protein that can handle various ionic strengths
	Question: Which protein can handle various ionic strengths?\\
	Answer: Encapsulin \\
	Visualisation: not needed\\
    \\\\
    Summary of this paper
\subsubsection{Methods}
    \begin{itemize}
        \item Materials used (Proteins, particles (100nm, 30 nm PS, DNA origami), Salt mixes, setup)
        \item chip prep (anti-fouling)
        \item encapsulin production (supplementary figure on encapsulin production)
        \item trapping at 0, 27 mM, 54 mM, 81 mM, 108 mM, 137 mM, going up and down with the concentration on 3 days/Gaps, ideally up and down on the same gap, for encapsulin, DNA and particles.
        \item analysis (steps and programs)
    \end{itemize}
\subsubsection{Results}
    Situation: trapping potential formula of particles and proteins\\
	Complication: ionic strength changes this\\
	Question: how does ionic strength change this?\\
	Answer: Calculation showing influence of ionic strength on potential for particles and proteins in theory (should be linear according to \cite{cametti2011Dielectric}?)\\
	Visualisation: not needed\\
    \\\\
    Situation: We understand the theory of ionic strength on trapping in theory\\
	Complication: we need to check in practice\\
	Question: How does ionic strength influence trapping probability in particles and DNA?\\
	Answer: show particle trapping probability/frequency at various ionic strength \\
	Visualisation: Figure 1: ionic strength vs trapping probability particles (100 nm, 30 nm PS, 40 nm Silica (?), DNA spheres)\\
    \\\\
    Situation: We know the trapping probability in correlation with the ionic strength\\
	Complication: the probability might change or stay the same, but the behaviour (duration, intensity) might change\\
	Question: Does the trapping behaviour change?\\
	Answer: Show trapping duration and intensity distribution of events (trapped and untrapped)\\
	Visualisation: Figure 2 Trapping behaviour of particles and DNA\\
    \\\\
    Situation: We see how trapping works with particles\\
    Complications: Proteins might behave differently\\
	Question: How do proteins react to high ionic strength?\\
	Answer: show protein probability/ frequency at various ionic strength \\
	Visualisation: Figure 3 ionic strength vs trapping probability particles (100 nm, 30 nm PS, 40 nm Silica (?)) ionic strength vs trapping probability proteins \\
    \\\\
    Situation: We know trapping probability of proteins\\
	Complication: their trapping duration and signal intensity could change as well \\
	Question: Does the trapping duration and signal strength change in correlation with the conductivity?\\
	Answer: Show trapping duration and intensity distribution of events (trapped and untrapped)\\
	Visualisation: Figure 4: show trapping behaviour of proteins \\
    \\\\
\subsubsection{Discussion}
What did go wrong? What are the limits of this data?\\
Compare behaviour of particles and DNA to theory.\\
Compare behaviour of proteins to theory.\\
Compare the behaviour of particles, DNA, and encapsulin. How are they different? How are they the same? \\
What do these results say about our setup? \\
What do these results say about the field?\\
\subsubsection{Outlook}
\FloatBarrier
